\begin{textsample}{POS Dim 3 – pe – Score 15.00 – pe/pe\_inf\_19.txt}  \label{ex:f3_pos_277}
2.7 ORGANIZAÇÃO DO TEMPO E ESPAÇO NA EDUCAÇÃO \textbf{INFANTIL} : A \textbf{ROTINA} \textbf{DIÁRIA} E O ESPAÇO FÍSICO A \textbf{rotina} na Educação \textbf{Infantil} é instrumento de reflexão sobre a prática pedagógica na busca de potencializar ações organizadas e flexibilizadas que atendam às necessidades de aprendizagem e desenvolvimento e promovam a autonomia das \textbf{crianças} nas interações e \textbf{brincadeiras}, no tempo e espaço.
Estudos realizados por Barbosa ( 2006 ) compreendem a \textbf{rotina} como uma categoria pedagógica da Educação \textbf{Infantil}, que opera como uma estrutura básica organizadora do cotidiano das \textbf{creches} ou pré-escola, sendo atividades recorrentes ou reiterativas na vida cotidiana coletiva, que são ressignificadas e planejadas de acordo com a ampliação das experiências das \textbf{crianças}.
Nessa dinâmica, o planejamento do professor ganha cada vez mais relevância no uso efetivo e qualificado do tempo ao arquitetar uma \textbf{rotina} flexível, promovendo reflexões metodológicas na utilização dos materiais disponíveis e \textbf{brincadeiras}, nas atividades individuais e em grupos, através de múltiplas linguagens, à medida que as ações da prática pedagógica apresentem \textbf{indicativos} para ampliação do conhecimento na trajetória de aprendizagem.
Quando se pensa em \textbf{rotina}, na \textbf{instituição} de Educação \textbf{Infantil}, relaciona -se também a ideia de espaço e tempo vividos.
No entanto, é importante a compreensão do que se trata esse tempo/espaço no currículo para constituição da formação humana.
É significativa a escolha de lugares para o desenvolvimento das \textbf{brincadeiras} e repouso, alimentação, \textbf{higiene}, atividades \textbf{lúdicas}, jogos e \textbf{brincadeiras}.
Deve -se considerar a importância de estimular nas \textbf{crianças} a realização de explorações, experimentações e descobertas que atuam no desenvolvimento de suas potencialidades físicas, cognitivas, motoras, afetivas e interativas, considerando suas características físicas, suas faixas \textbf{etárias} e sua cultura.
O tempo e o espaço também devem ser organizados intencionalmente de forma a promover o desenvolvimento integral, o que remete à importância da sensibilidade do professor diante dos materiais e móveis que ocupam esses espaços e sua relevância significativa em contribuir para aprendizagem e desenvolvimento desses sujeitos de direitos, desejantes, ativos, cognoscentes.
A valorização do desenvolvimento integral, nos aspectos cognitivo, físico e socioemocional, assim como a organização do espaço e do tempo nas \textbf{creches} e \textbf{pré-escolas} tornam -se substanciais quando se apresentam comprometidas com os direitos de aprendizagem e de desenvolvimento, pois “ a organização do espaço se constitui em um \textbf{parceiro} pedagógico de excelência.
Quanto mais rico e desafiador este espaço for, mais qualificadas serão as aprendizagens das \textbf{crianças} ” ( HORN, 2004 ).
Portanto, as \textbf{instituições} de Educação \textbf{Infantil} devem garantir espaços físicos limpos, seguros, inclusivos, acolhedores e desafiadores, com acessibilidade, estética, ventilação, insolação, luminosidade, acústica, com \textbf{higiene} e interatividade, de sorte a permitir a participação efetiva nas explorações e descobertas nas relações e interações entre crianças/crianças ; crianças/professor ; Como explica Rossetti-Ferreira et al ( 2007 ), ao citar que os espaços devem ser sempre atraentes e estimulantes na Educação \textbf{Infantil}, o professor necessita estar \textbf{atento} para as observações e mudanças, de modo a acompanhar o desenvolvimento dos \textbf{bebês} e das \textbf{crianças} e seus interesses por coisas novas.
Nessa perspectiva, a \textbf{rotina} e o espaço físico nas \textbf{creches} e \textbf{pré-escolas} estão imbricados e devem estar comprometidos com a aprendizagem e desenvolvimento na Educação \textbf{Infantil}, pois reconhecer que os \textbf{bebês} e as \textbf{crianças} estão sempre aprendendo, é imprescindível!

% matched lemmas: atento, bebê, brincadeira, creche, criança, diário, etário, higiene, indicativo, infantil, instituição, lúdico, parceiro, pré-escola, rotina
\end{textsample}
